\documentclass[10pt,a4paper]{amsart}
\usepackage[margin=19mm]{geometry}
\usepackage{amsmath,amssymb,mathtools}
\usepackage[scr=rsfs,cal=euler]{mathalfa}
\usepackage{xfrac}
\usepackage[unicode,hidelinks,bookmarks=false]{hyperref}
\newcommand{\ie}{i.e.\ }
\newcommand{\cf}{cf.\ }
\newcommand{\resp}{respectively\ }
\newcommand{\Subsection}{\subsection}
\providecommand{\nobreakdash}{\nobreak\hbox{-}}
\setlength{\emergencystretch}{2em}
\multlinegap0pt
\usepackage{amssymb}
\newtheorem{theorem}{Theorem}
\title{Cartan meets Cramér-Rao}
\author{Sunder Ram Krishnan}
\date{Source: arXiv:2511.15612v2 (CC BY 4.0)}
\begin{document}
\maketitle

\noindent\emph{Source and licence.} Cartan meets Cramér-Rao. Version arXiv:2511.15612v2 (CC BY 4.0), distributed by arXiv under the Creative Commons Attribution 4.0 International licence, http://creativecommons.org/licenses/by/4.0/, as recorded in the arXiv licence metadata for this exact version and shown on the abstract page https://arxiv.org/abs/2511.15612v2. Full licence notice: \texttt{licenses/arxiv-2511.15612-CC-BY-4.0.txt}.\par\smallskip

\noindent\textit{Editorial scope.} Counted unit: the article's theorem thm:equivalence (source0.tex lines 679--691). The article's complete original proof of it is delivered with it (source0.tex lines 693--713, one \texttt{proof} environment of 21 lines). No mathematics is written, repaired or re-derived: the statement and the proof are the article's own lines, in the article's own notation, and no retained line is substituted or re-flowed.  Retained uncounted as the source-local context the proof needs: lines 662--664.  Nothing was omitted from any retained span, so the package carries no omission record.  No retained line contains a citation, so the wrapper carries no bibliography and no bibliography entry.  No retained line contains a figure environment, an image-inclusion command or drawing code, so no image is delivered and the package declares no asset dependency.  Editorial formatting only, in addition: an AMS \texttt{amsart} preamble that restates the article's own declarations and macros used by the retained lines, namely the environments \texttt{theorem} on separate counters, together with the standard packages those lines need.  The retained environments print in this document as Theorem 1 in order of appearance, whereas the article prints the same items with the numbering of its own full document.  The complete native source of the article as distributed in the arXiv e-print for this exact version is bundled unchanged as \texttt{sources/189461/source0.tex}.
\begin{equation}\label{eq:proj-cond}
(T(X)-\theta)s_\theta \in T_m(\theta).
\end{equation}

\begin{theorem}[Cartan viewpoint of $m$-th order efficiency]
\label{thm:equivalence}
For any $m\ge1$, the following are equivalent:
\begin{enumerate}
\item The projection condition \eqref{eq:proj-cond} holds for all $\theta$.
\item The prolonged section $j^m s(\theta)$ takes values in $\Sigma_m$ and
the total derivative $D_\theta^{(m)}$ restricted to $\Sigma_m$ is tangent to
$\Sigma_m$; i.e.
\[
D_\theta^{(m)}\big|_{j^ms(\theta)}\in T_{j^m s(\theta)}\Sigma_m.
\]
\end{enumerate}
\end{theorem}

\begin{proof}
In coordinates on $J^m(E)$,
\[
D_\theta^{(m)}=
\partial_\theta +
s_1\partial_s +
\cdots +
s_m\partial_{s_{m-1}}.
\]
The constraint defining $\Sigma_m$ is linear in $(s,s_1,\dots,s_m)$.  A
vector is tangent to $\Sigma_m$ iff it annihilates the constraint
one--form
\[
\alpha_m
=
(T-\theta)sd\theta - c_{1,\theta}ds - \cdots - c_{m,\theta}ds_{m-1}.
\]
Evaluating $\alpha_m$ on $D_\theta^{(m)}$ gives exactly the defect of
$(T-\theta)s_\theta$ from the span of $s_\theta^{(1)},\dots,s_\theta^{(m)}$.
Thus $\alpha_m(D_\theta^{(m)})=0$ iff the projection condition holds.
\end{proof}

\end{document}
